problem:  
Let \((M^{2n}, g, J)\) be a compact **Kähler manifold** with positive scalar curvature \(\mathrm{Scal}_M>0\). For each point \(p\in M\), define  
\[
H_{\max}(p) := \max_{|X|=1} H(X), \qquad H_{\min}(p) := \min_{|X|=1} H(X),
\]
where \(H(X)=\frac{R(X,\bar{X},X,\bar{X})}{|X|^4}\) is the holomorphic sectional curvature for \(X\in T^{1,0}_p M\setminus\{0\}\).  
Set
\[
 \mathrm{osc}_H(p) := H_{\max}(p)-H_{\min}(p),
\qquad 
\kappa(M):=\frac{\int_M \mathrm{osc}_H(p)^2\, dV_g}{\int_M \mathrm{Scal}_M^2\, dV_g}.
\]
Thus \(\kappa(M)\) measures the pointwise variation of holomorphic sectional curvature relative to the positive scalar curvature.  

Let \((N,h)\) be a compact Riemannian manifold with nonpositive sectional curvature, and \(u:(M,g)\to(N,h)\) a harmonic map minimizing the energy functional.  

**Question:**  
Does there exist a constant \(C_n>0\), depending only on the complex dimension \(n\), such that if \(\mathrm{Scal}_M>0\) and \(\kappa(M)< C_n\), then every harmonic map \(u:(M,g)\to(N,h)\) is constant?  
Equivalently, does *sufficiently small variation of holomorphic sectional curvature*—that is, a nearly uniform holomorphic curvature pinching—together with positive scalar curvature enforce harmonic rigidity in the scalar-positive, Kähler setting?

Why is it a "good" problem:  
This version resolves the earlier ambiguities and introduces a **natural pinching invariant** \(\kappa(M)\) whose geometric meaning is clear: it measures the dispersion of holomorphic sectional curvature values across tangent directions at each point.  Small \(\kappa(M)\) is equivalent to requiring that \(M\) be close (in curvature sense) to a complex space form, which includes Kähler–Einstein manifolds of constant holomorphic sectional curvature as extrema where Eells–Sampson rigidity is known.  

The question explores whether this *nearness to constant holomorphic curvature* can substitute for positive Ricci curvature in guaranteeing harmonic rigidity — an uncharted analytical frontier that connects **Kähler curvature pinching**, **Bochner methods**, and **harmonic map theory**.  A positive resolution would establish a quantitative bridge between curvature uniformity and analytic rigidity; a counterexample would reveal precise limits of curvature control in the scalar-positive Kähler regime.  The problem is conceptually simple, geometrically natural, open, and it demands cross-disciplinary insight—meeting the criteria for a truly “good” research-level problem.